% Complete definitions supporting the compact main-paper method.
% Proposed V4 recipe; biological supervision is retained only as an ablation.
\section{Detailed architecture and training objectives}
\label{app:paper-specification}
\label{sec:v4-method}

\subsection{Notation and computational structure}
\label{sec:v4-problem}
We provide the complete definitions underlying Section~\ref{sec:paper-v4-method}.
We use the same symbols as the main text: bold letters denote vectors or
matrices, and plain indexed letters their scalar entries.
Let $\vR=\{r_1,\ldots,r_{N_R}\}$ and
$\vE=\{e_1,\ldots,e_{N_E}\}$ denote the unique training endpoints, with
association graph $\vD\subseteq\vR\times\vE$ and binary adjacency
$Y^{\mathrm{tr}}_{ij}=\mathbb{1}[(r_i,e_j)\in\vD]$.
Figure~\ref{fig:v4-architecture} summarizes the independently computable
representations. Their score measures retrieval compatibility, rather than
a calibrated probability of catalytic activity. Both training stages use
enzyme--reaction associations; added annotation losses are evaluated only
in the ablation in Section~\ref{sec:v4-biology}.

Throughout, $\vunit{\mathbf{x}}=\mathbf{x}/\max(\lVert\mathbf{x}\rVert_2,\epsilon)$ denotes numerical
$\ell_2$ normalization, $\vpos{x}=\max(0,x)$, $[\mathbf{x};\mathbf{y}]$ denotes concatenation,
and $\vLN$ denotes layer normalization. All attention distributions exclude
padding. Networks with different labels have independent parameters unless
sharing is explicitly stated.

\begin{figure}[tb]
\centering
\begin{tikzpicture}[
    box/.style={draw=black!55,rounded corners,align=center,font=\small,
                text width=4.65cm,minimum height=1.0cm,inner sep=5pt},
    arrow/.style={-{Latex[length=2mm]},thick},node distance=0.55cm]
  \node[box,fill=blue!5] (protein) {Protein sequence\\Frozen ProtT5\\residue features};
  \node[box,fill=orange!7,right=1.2cm of protein] (reaction) {Reaction input\\Frozen model features\\and chemistry};
  \node[box,fill=blue!5,below=of protein] (pe) {Global + SLEEC + $K$ slots\\Feature-wise fusion\\512-D protein vector};
  \node[box,fill=orange!7,below=of reaction] (re) {Molecular side pooling\\Modality attention\\512-D reaction vector};
  \node[box,below=of pe] (ph) {Protein residual refinement\\512-D neural representation};
  \node[box,below=of re] (rh) {Reaction residual refinement\\512-D neural representation};
  \node[box,fill=black!3,below=of ph] (pa) {Append dictionary block\\Support from 32\\training protein neighbors};
  \node[box,fill=black!3,below=of rh] (ra) {Append dictionary block\\Similarity to all\\training reactions};
  \node[draw=black!55,rounded corners,align=center,font=\small,
        text width=9.6cm,below=0.65cm of $(pa.south)!0.5!(ra.south)$] (score)
       {Dot product: $S(r,e)=0.6\,\mathbf{u}_r^\top \mathbf{u}_e+0.4\,\mathbf{a}_r^\top \mathbf{a}_e$\\
        Independent vectors for both $R\!\to\!E$ and $E\!\to\!R$};
  \draw[arrow] (protein)--(pe); \draw[arrow] (reaction)--(re);
  \draw[arrow] (pe)--(ph); \draw[arrow] (re)--(rh);
  \draw[arrow] (ph)--(pa); \draw[arrow] (rh)--(ra);
  \draw[arrow,dashed,draw=blue!60!black] (protein.west) -- ++(-0.4,0) |- (pa.west);
  \draw[arrow,dashed,draw=orange!70!black] (reaction.east) -- ++(0.4,0) |- (ra.east);
  \draw[arrow] (pa.south)--(score.north -| pa.south);
  \draw[arrow] (ra.south)--(score.north -| ra.south);
\end{tikzpicture}
\caption{Inference structure of \methodname, with $K=4$ by default. Dashed paths supply frozen raw
features to the dictionary maps. Dictionary coordinates are computed
from frozen raw features and training associations, independently of the neural
refinement heads. The arrows into the appended representations indicate
composition, not that dictionary features are generated from the heads.
The residual heads are trained with association cross-entropy and an
identity penalty. The protein fused-view multiplier
and refinement multiplier are specified in Section~\ref{sec:v4-inference}.}
\label{fig:v4-architecture}
\end{figure}


\subsection{Protein encoder: global, SLEEC, and learned residue views}
\label{sec:v4-protein}
For an enzyme sequence $e$ with $L$ retained residues, a frozen ProtT5 model
provides contextual residue vectors
$\mathbf{H}_e=[\mathbf{h}_1,\ldots,\mathbf{h}_L]^\top\in\mathbb{R}^{L\times1024}$
\citep{elnaggar2022prottrans}. Here $\mathcal V_e=\{1,\ldots,L\}$ indexes
valid residues; $e$ is suppressed on per-residue vectors and attention weights.
We construct $K+2$ protein views: one global
view, one view defined by a frozen SLEEC functional-residue scorer
\citep{gong2026sleec}, and $K$ learned residue views, with $K=4$ by default. The scorer is a
pretrained resource; it is not fitted to each retrieval target in phase 1.

\paragraph{Global and SLEEC views.}
The global view is the mean of the retained residue embeddings,
\begin{equation}
  \mathbf{x}_e^{P}=\frac{1}{L}\sum_{i=1}^{L}\mathbf{h}_i.
  \label{eq:v4-global}
\end{equation}
The frozen scorer computes a scalar logit $\ell_i$ for each residue using a
$1024\!\to\!256\!\to\!1$ MLP with a ReLU hidden activation. We center these
logits within the protein, clip them, and form a smoothed attention distribution:
\begin{align}
  \bar\ell &= L^{-1}\sum_{i=1}^{L}\ell_i,
  & \tilde\ell_i &= \operatorname{clip}(\ell_i-\bar\ell,-10,10),\\
  p_i^{\mathrm{S}} &= \frac{\exp(\tilde\ell_i)}{\sum_{j=1}^{L}\exp(\tilde\ell_j)},
  & a_i^{\mathrm{S}} &= (1-\eta)p_i^{\mathrm{S}}+\frac{\eta}{L},
  \quad \eta=0.05,\\
  \mathbf{x}_e^{S} &= \sum_{i=1}^{L}a_i^{\mathrm{S}}\mathbf{h}_i.
  \label{eq:v4-sleec}
\end{align}
The vector $\mathbf{x}_e^{S}\in\mathbb R^{1024}$ is a protein-level summary:
it averages the residue embeddings $\mathbf h_i$ with nonnegative weights
$a_i^{\mathrm S}$ that sum to one. The superscripts $P$ and $S$ identify
the global ProtT5 and SLEEC-weighted features, respectively; superscripts
$k\in\{1,\ldots,K\}$ identify the learned features. The threshold used for
SLEEC diagnostics does not select or remove residues in this encoder.

\paragraph{Learned residue views.}
We augment the global and SLEEC summaries with $K$ attention-pooled views
learned from enzyme--reaction associations. These views provide adaptive
pooling of contextual residue features before feature-wise fusion.

A shared residue adapter maps each contextual residue to $d_h=256$ dimensions:
\begin{equation}
  \tilde{\mathbf{h}}_i=\vGELU(\mathbf{W}_{\mathrm{res}}\vLN(\mathbf{h}_i)+\mathbf{b}_{\mathrm{res}}),\qquad
  \mathbf{k}_i=\vunit{\mathbf{W}_{\mathrm{key}}\tilde{\mathbf{h}}_i},\qquad
  \mathbf{v}_i=\mathbf{W}_{\mathrm{val}}\tilde{\mathbf{h}}_i+\mathbf{b}_{\mathrm{val}}.
  \label{eq:v4-residue-adapter}
\end{equation}
The unit-normalized key $\mathbf{k}_i$ determines residue $i$'s similarity to a query;
the value $\mathbf{v}_i$ contains the features aggregated into the view. The key
projection $\mathbf{W}_{\mathrm{key}}$ is bias-free, while the value projection is affine.
Queries $\mathbf{q}_k\in\mathbb{R}^{256}$ are orthogonally initialized model parameters
shared across enzymes and optimized during phase 1. Normalizing each query
gives $\hat{\mathbf{q}}_k=\vunit{\mathbf{q}_k}$. Query $k$ produces raw view
$\mathbf{x}_e^{k}$:
\begin{align}
  \gamma &=10\vsigmoid(\rho_\gamma),\qquad \gamma_{\mathrm{init}}=4,\\
  p_{ki}&=\frac{\exp\!\left(\gamma\,\hat{\mathbf{q}}_k^{\top}\mathbf{k}_i\right)}
                     {\sum_{j=1}^{L}\exp\!\left(\gamma\,\hat{\mathbf{q}}_k^{\top}\mathbf{k}_j\right)},
  & a_{ki}&=(1-\eta)p_{ki}+\frac{\eta}{L},\\
  \mathbf{x}_e^{k}&=\sum_{i=1}^{L}a_{ki}\mathbf{v}_i,
  & k&\in\{1,\ldots,K\}.
  \label{eq:v4-slots}
\end{align}
For each query, the softmax is taken over valid residues, so $p_{ki}\geq0$
and $\sum_i p_{ki}=1$. The learned scalar $\rho_\gamma$ parameterizes a
bounded scale $\gamma\in(0,10)$: larger values concentrate attention on
residues with higher query--key cosine similarity. The uniform mixture gives
$a_{ki}\geq\eta/L$ and $\sum_i a_{ki}=1$. With $\eta=0.05$, the resulting
view is a mixture of the attention-weighted mean (95\%) and the uniform
mean (5\%) of the same value vectors. These weights quantify feature
aggregation, not probabilities that residues are catalytic.

Thus, $\mathbf{q}_k$ defines a shared pooling function, while $a_{ki}$ and
$\mathbf{x}_e^{k}$ are enzyme-specific. The phase-1 objective updates the queries,
adapter, and projections through the fused enzyme embedding; entropy and
diversity penalties discourage overly concentrated or redundant attention
(Section~\ref{sec:v4-phase1}). Views can pool distant residues and overlap,
without assigned biochemical roles. Because pooling and fusion use only
enzyme features, the views and their fused representation can be cached
once the encoder is frozen and reused across candidate reactions.
SLEEC contributes a separate view, not a bias to the learned attention logits.

\paragraph{Choice and interpretation of the view count.}
The view count $K$ sets the number of adaptive summaries available before
fusion. With $K=1$, all learned residue information passes through one
summary; larger $K$ can retain complementary patterns, but also increases
query, projection, and gate parameters. Views may overlap, so a larger count
does not guarantee additional useful information. We retain the original
$K=4$ default as a compromise across screening and bidirectional retrieval:
the completed $K\in\{1,2,4,8\}$ ablation finds metric-dependent trade-offs
rather than a universally best count (Appendix~\ref{app:v4-k-ablation}).
This single-seed study does not establish an optimal capacity or imply that
four views correspond to four functional regions or biochemical properties.
The residue queries remain frozen throughout phase-2 refinement.

\paragraph{Feature-wise gated fusion.}
The extracted enzyme features are combined by feature-wise gated fusion.
Each feature is independently projected into a common 256-dimensional space
using a network of the form
\begin{equation}
  F_{a\to b}(\mathbf{x})=\mathbf{W}_2\vDrop\!\left(\vGELU(\mathbf{W}_1\vLN(\mathbf{x})+\mathbf{b}_1)\right)+\mathbf{b}_2,
  \label{eq:v4-protein-ffn}
\end{equation}
whose hidden width is 512 and whose training dropout probability is 0.1.
The input width is $a=1024$ for views $P,S$ and $a=256$ for learned views.
Define $\mathbf{y}_e^{v}=\vunit{F_v(\mathbf{x}_e^{v})}$ for
$v\in\mathcal I=\{P,S,1,\ldots,K\}$, and let $V=|\mathcal I|=K+2$.
A $256V\!\to\!256\!\to\!256V$ gate MLP with GELU maps their concatenation
to logits $\mathbf{G}_e\in\mathbb{R}^{V\times256}$, with rows ordered as
$P,S,1,\ldots,K$. The gate softmax is taken over
views \emph{separately for each feature coordinate} $d=1,\ldots,256$:
\begin{align}
  \omega_{e,vd}
    &=(1-\zeta)\frac{\exp(G_{e,vd})}{\sum_{u\in\mathcal I}\exp(G_{e,ud})}
       +\frac{\zeta}{K+2},\qquad \zeta=0.05,\\
  \bar y_{e,d}&=\sum_{v\in\mathcal I}\omega_{e,vd}\,y_{e,d}^{v}.
  \label{eq:v4-view-gate}
\end{align}
For each enzyme and coordinate, $\sum_v\omega_{e,vd}=1$: each fused coordinate
is a weighted average of that coordinate across the projected views.
The weights vary with both the enzyme and the coordinate, allowing different
views to supply different features. The final gate layer is initialized to
zero, giving equal initial view weights. The floor is $\zeta/(K+2)$ per view
and coordinate, hence $\zeta/6$ at the default $K=4$.

Two additional networks of the form in Eq.~\eqref{eq:v4-protein-ffn} produce
separately normalized global and fused outputs in $\mathbb{R}^{512}$:
\begin{align}
  \mathbf{g}_e &= \vunit{F_g(\mathbf{x}_e^{P})}, & \mathbf{f}_e&=\vunit{F_f(\bar{\mathbf{y}}_e)},\\
  s&=0.5\vsigmoid(\rho_s), & \mathbf{b}_e&=\vunit{\mathbf{g}_e+s \mathbf{f}_e}.
  \label{eq:v4-protein-output}
\end{align}
Here $\mathbf g_e$ is computed directly from the sequence mean, while
$\mathbf f_e$ is computed from the combined views. Their separate projections
both output 512-dimensional vectors, making the final addition well-defined.
The scalar $s$ is shared across proteins, is initialized to 0.1, and remains
strictly between 0 and 0.5. It controls the strength of the entire fused branch;
the weights $\omega_{e,vd}$ instead control the contributions of individual
views within that branch. The direct global branch retains sequence context
alongside the information supplied by fusion.

\paragraph{Sequence order and interpretation.}
No additional positional encoding is applied in the pooling module. Its
queries attend to the contextual content of ProtT5 vectors. Permuting already
computed residue vectors and their masks together leaves this pooling operation
unchanged, whereas changing the amino-acid sequence can change the contextual
vectors. Neither a slot nor a coordinate of the final 512-dimensional vector
is assigned a named catalytic function. Attention weights are learned
aggregation weights, not validated active-site labels.

\subsection{Reaction encoder: molecular sides and modality fusion}
\label{sec:v4-reaction}
We use four feature sources: a 768-dimensional ReactionT5v2 reaction embedding
\citep{sagawa2023reactiont5,sagawaReactionT5v2}, 768-dimensional Uni-Mol2
molecule embeddings \citep{ji2024unimol2}, 256-dimensional ChIRo molecule
embeddings \citep{adams2022chiro}, and a 617-dimensional fixed chemistry
vector defined in Appendix~\ref{app:v4-chemistry}. Its four blocks summarize
participant substructures, molecular properties, set statistics, and cofactor
presence. Their feature extractors are frozen; molecular pooling, modality
projections, and fusion are trainable.

\paragraph{Side-specific molecule pooling.}
For one reaction (with index $r$ suppressed), molecular modality
$m\in\{\mathrm{UniMol2},\mathrm{ChIRo}\}$, and side
$s\in\{-,+\}$, let $\mathbf{x}_{s,j}^{m}$ be the embedding of molecule $j$.
Independent linear attention scorers for the two sides produce
\begin{equation}
  \pi_{s,j}^{m}=
    \frac{\exp((\mathbf{w}_s^m)^\top \mathbf{x}_{s,j}^{m}+b_s^m)}
         {\sum_t\exp((\mathbf{w}_s^m)^\top \mathbf{x}_{s,t}^{m}+b_s^m)},
  \qquad
  \mathbf{p}_s^m=\sum_j\pi_{s,j}^{m}\mathbf{x}_{s,j}^{m}.
  \label{eq:v4-molecule-pool}
\end{equation}
The configured \texttt{directional\_delta} composition is
\begin{equation}
  \mathbf{c}_r^m=[\mathbf{p}_-^m;\ \mathbf{p}_+^m;\ \mathbf{p}_+^m-\mathbf{p}_-^m;\ |\mathbf{p}_+^m-\mathbf{p}_-^m|].
  \label{eq:v4-directional}
\end{equation}
It has dimension 3072 for Uni-Mol2 and 1024 for ChIRo. Direction is represented
by the ordered sides and signed difference when physical sides are supplied.
The first two blocks preserve the pooled molecular context of each side;
the last two expose the sign and magnitude of coordinate-wise contrasts.
These contrasts are between learned summaries, rather than an explicit
atom-mapped list of bond changes. Only Uni-Mol2 and ChIRo use this pooling:
ReactionT5v2 and the fixed chemistry descriptor already provide one vector
per reaction.

\paragraph{Input semantics.}
For EnzymeMap, inputs retain their physical reactant and product sides.
ReactZyme's released participant-set inputs use the established self-reaction
adapter, placing the same complete set on both sides. The two side poolers
still have different parameters, so their pooled vectors need not be equal.
Their difference on these self-reactions must not be interpreted as an
observed chemical transformation direction. The choice of chemical input
representation is separate from the choice of retrieval direction.

\paragraph{Modality tokens and attention.}
For $m\in\{\mathrm{T5},\mathrm{UniMol2},\mathrm{ChIRo},\mathrm{chem}\}$,
let $\mathbf{c}_r^m$ denote the corresponding reaction input. Each independent modality
MLP has widths $d_m\!\to\!1024\!\to\!512\!\to\!512$, with ReLU after
the first two linear layers. A shared layer-normalization module then gives
the modality token
\begin{equation}
  \mathbf{t}_r^m=\vLN(\operatorname{MLP}_m(\mathbf{c}_r^m))\in\mathbb{R}^{512}.
\end{equation}
An attention MLP of widths $512\!\to\!512\!\to\!1$, with a ReLU hidden
activation, scores the tokens. If $\mathcal M_r$ is the set of available
modalities,
\begin{equation}
  \alpha_r^m=
  \frac{\exp(A(\mathbf{t}_r^m))}{\sum_{n\in\mathcal M_r}\exp(A(\mathbf{t}_r^n))},
  \quad
  \mathbf{t}_r=\sum_{m\in\mathcal M_r}\alpha_r^m \mathbf{t}_r^m,
  \quad
  \mathbf{b}_r=\vunit{F_{\mathrm{out}}(\mathbf{t}_r)}.
  \label{eq:v4-reaction-output}
\end{equation}
The output MLP $F_{\mathrm{out}}$ has widths $512\!\to\!1024\!\to\!512$ and a ReLU
hidden activation. Unavailable modalities are masked out of attention.
Molecule attention therefore combines molecules within one modality and
side, whereas modality attention combines the four reaction-level tokens.
Reaction modality dropout and projection dropout are zero in the reported
recipe. The resulting $\mathbf{b}_r$ and $\mathbf{b}_e$ inhabit a shared 512-dimensional
retrieval space.

\subsection{Phase 1: anchor-balanced, decoupled multi-positive alignment}
\label{sec:v4-phase1}
Training minibatches sample association rows and deduplicate their endpoint
identifiers before constructing a reaction-by-enzyme score matrix. Let
$\mathcal R_B$ and $\mathcal E_B$ be these distinct endpoints. Positives are
identified from \emph{all known training associations between endpoints in
the batch}, rather than only the association rows sampled for that batch.
For a reaction anchor $r$, define
\begin{equation}
  P_B(r)=\{e\in\mathcal E_B:(r,e)\in\vD\},\qquad
  N_B(r)=\mathcal E_B\setminus P_B(r),
\end{equation}
and define $P_B(e)$ and $N_B(e)$ analogously for enzyme anchors. Only anchors
with at least one positive and one remaining candidate enter the respective
directional average.

The logits are $s_{re}^{(1)}=\mathbf{b}_r^\top \mathbf{b}_e/\tau_1$ with $\tau_1=0.2$. For each known positive $p$ of an anchor $a$, the
decoupled loss is
\begin{align}
  \ell_{\mathrm{dec}}(a,p)
   &=-\log\frac{\exp(s_{ap}^{(1)})}
       {\exp(s_{ap}^{(1)})+\sum_{n\in N_B(a)}\exp(s_{an}^{(1)})}\\
   &=\log\left(1+\sum_{n\in N_B(a)}\exp(s_{an}^{(1)}-s_{ap}^{(1)})\right).
   \label{eq:v4-decoupled}
\end{align}
Other known positives are excluded from the denominator for this individual
positive. Thus every positive must compete with the remaining candidates,
without competing with its known positive partners. This differs from a
single log-sum over all positives, which can be dominated by one easy positive.

Let $\mathcal A_R$ and $\mathcal A_E$ be the valid reaction and enzyme anchors
in the batch. We first average over positives within an anchor, then average
over anchors, and finally average the two directions:
\begin{align}
  \mathcal L_{R\to E}^{(1)}
    &=\frac{1}{|\mathcal A_R|}\sum_{r\in\mathcal A_R}
      \frac{1}{|P_B(r)|}\sum_{e\in P_B(r)}\ell_{\mathrm{dec}}(r,e),\\
  \mathcal L_{E\to R}^{(1)}
    &=\frac{1}{|\mathcal A_E|}\sum_{e\in\mathcal A_E}
      \frac{1}{|P_B(e)|}\sum_{r\in P_B(e)}\ell_{\mathrm{dec}}(e,r),\\
  \mathcal L_{\mathrm{align}}^{(1)}
    &=\tfrac12\mathcal L_{R\to E}^{(1)}+
      \tfrac12\mathcal L_{E\to R}^{(1)}.
  \label{eq:v4-anchor-balanced}
\end{align}
A direction with no valid anchor contributes zero. The normalization prevents
an anchor's contribution from growing simply in proportion to its number of
positives in the matrix. It does not make anchors equally frequent over an
epoch when association rows are sampled. Unobserved pairs are still used as
contrastive negatives with weight 1; they are not measured inactive pairs.

\paragraph{Attention regularization.}
The protein encoder also returns two small penalties on the \emph{raw}
learned-slot distributions $p_{ki}$, before the uniform mixture in
Eq.~\eqref{eq:v4-slots}. For a protein of length $L$, define
\begin{align}
  H_k&=-\sum_i p_{ki}\log p_{ki},\qquad
  H_{\min}(L)=\min\{\tfrac12\log L,\log4\},\\
  E(e)&=\frac1K\sum_{k=1}^{K}\vpos{H_{\min}(L)-H_k}^{2}.
  \label{eq:v4-entropy}
\end{align}
This penalty discourages attention with very small effective support. To
discourage nearly identical attention patterns, center each raw distribution
against uniform attention: $\bar{\mathbf{p}}_k=\mathbf{p}_k-L^{-1}\mathbf 1$,
where $\mathbf{p}_k=(p_{ki})_{i\in\mathcal V_e}$. For eligible
unordered pairs $\mathcal Q_e=\{(k,l):k<l,\ \|\bar{\mathbf{p}}_k\|_2>10^{-6},
\ \|\bar{\mathbf{p}}_l\|_2>10^{-6},\ L>1\}$, let
\begin{equation}
  D(e)=\frac{1}{\max(1,|\mathcal Q_e|)}
     \sum_{(k,l)\in\mathcal Q_e}
       \vpos{\frac{\bar{\mathbf{p}}_k^\top\bar{\mathbf{p}}_l}
                         {\|\bar{\mathbf{p}}_k\|_2\|\bar{\mathbf{p}}_l\|_2}-0.9}^{2}.
  \label{eq:v4-diversity}
\end{equation}
The main-text penalties are $\mathcal L_{\mathrm{ent}}=|\mathcal E_B|^{-1}
\sum_{e\in\mathcal E_B}E(e)$ and $\mathcal L_{\mathrm{div}}=|\mathcal E_B|^{-1}
\sum_{e\in\mathcal E_B}D(e)$. Uniform heads have no meaningful centered direction and are excluded from
this comparison. The complete phase-1 objective is
\begin{equation}
  \boxed{\mathcal L^{(1)}=
      \mathcal L_{\mathrm{align}}^{(1)}
      +0.01\,\frac{1}{|\mathcal E_B|}\sum_{e\in\mathcal E_B}E(e)
      +0.01\,\frac{1}{|\mathcal E_B|}\sum_{e\in\mathcal E_B}D(e).}
  \label{eq:v4-phase1-total}
\end{equation}
The safeguards in Eqs.~\eqref{eq:v4-entropy}
and~\eqref{eq:v4-diversity} encourage numerical and representational diversity;
they do not establish a biological interpretation of the attention heads.

\subsection{Phase 2: full-graph residual refinement}
\label{sec:v4-phase2}
After retaining the epoch-10 phase-1 snapshot, we freeze both F3 encoders and
cache their unadjusted outputs $\mathbf{b}_r,\mathbf{b}_e$ for the unique training endpoints.
Two independent residual heads $h_R,h_E$ each comprise layer normalization,
a $512\!\to\!1024$ linear layer, GELU, and a $1024\!\to\!512$ linear
layer. The last layer's weights and biases are initialized to zero. During
phase-2 training,
\begin{equation}
  \mathbf{u}_r^{\mathrm{tr}}=\vunit{\mathbf{b}_r+\kappa h_R(\mathbf{b}_r)},\qquad
  \mathbf{u}_e^{\mathrm{tr}}=\vunit{\mathbf{b}_e+\kappa h_E(\mathbf{b}_e)},\qquad \kappa=0.2.
  \label{eq:v4-residual-training}
\end{equation}
Consequently, both heads initially preserve the normalized base
representations. No gradient is propagated to the frozen F3 encoders.

\paragraph{Uniform-positive cross-entropy over the complete graph.}
We abbreviate $s_{ij}^{(2)}=s_{r_i e_j}^{(2)}=(\mathbf{u}_{r_i}^{\mathrm{tr}})^\top
\mathbf{u}_{e_j}^{\mathrm{tr}}/\tau_2$, with $\tau_2=0.2$. Let
$d_i^R=\sum_jY^{\mathrm{tr}}_{ij}$ and $d_j^E=\sum_iY^{\mathrm{tr}}_{ij}$. Every retained training
endpoint has at least one association. The directional objectives are
\begin{align}
  \mathcal L_{R\to E}^{(2)}
    &=\frac{1}{N_R}\sum_i\left[
       \log\sum_{j=1}^{N_E}\exp(s_{ij}^{(2)})
       -\frac{1}{d_i^R}\sum_jY^{\mathrm{tr}}_{ij}s_{ij}^{(2)}\right],\\
  \mathcal L_{E\to R}^{(2)}
    &=\frac{1}{N_E}\sum_j\left[
       \log\sum_{i=1}^{N_R}\exp(s_{ij}^{(2)})
       -\frac{1}{d_j^E}\sum_iY^{\mathrm{tr}}_{ij}s_{ij}^{(2)}\right],\\
  \mathcal L_{\mathrm{CE}}^{(2)}
    &=\tfrac12\mathcal L_{R\to E}^{(2)}+
      \tfrac12\mathcal L_{E\to R}^{(2)}.
  \label{eq:v4-fullgraph}
\end{align}
These are cross-entropies with a uniform target distribution over each
anchor's known positives. Unlike phase 1, the denominator includes all
opposite-type training endpoints, including all positives. Phase 2 therefore
changes both the coverage of the contrastive bank and the form of the loss;
it is not simply Eq.~\eqref{eq:v4-decoupled} evaluated on a larger batch.

\paragraph{Preserving the original geometry.}
An identity penalty limits changes to the frozen base representations:
\begin{equation}
  \mathcal L_{\mathrm{id}}=
  \frac{1}{2N_R}\sum_i\left[1-(\mathbf{u}_{r_i}^{\mathrm{tr}})^\top\vunit{\mathbf{b}_{r_i}}\right]
  +\frac{1}{2N_E}\sum_j\left[1-(\mathbf{u}_{e_j}^{\mathrm{tr}})^\top\vunit{\mathbf{b}_{e_j}}\right].
  \label{eq:v4-identity}
\end{equation}
The complete refinement objective is
\begin{equation}
  \boxed{\mathcal L^{(2)}=
     \mathcal L_{\mathrm{CE}}^{(2)}+10\mathcal L_{\mathrm{id}}.}
  \label{eq:v4-phase2-total}
\end{equation}
Both heads are optimized jointly for 100 updates, each of which includes
the entire compact training graph. Added biological-label losses are not
part of this objective.

\subsection{Biological-supervision ablation}
\label{sec:v4-biology}
This ablation is excluded from the proposed method. We retain its
definition to document the mixed experimental outcomes. The ablated
variant, \vfourBio, adds supervision to the same refinement heads.
It constrains relationships among reaction representations and, separately,
among enzyme representations. The association loss connects the two endpoint
types. No classification heads, biological feature coordinates, additional
residue queries, or inference-time annotation inputs are introduced.

\paragraph{Annotation families and confidence.}
We use three families $\mathcal F=\vFamilies$. Each training endpoint can
belong to multiple categories. A membership in category $c$ has confidence
$w_{ic}\in(0,1]$; missing memberships create no annotation evidence.
For EC, each specified prefix is retained. A category at hierarchy depth
$d\in\{1,2,3,4\}$ has weight
$\xi_c\in\{0.125,0.25,0.5,1\}$, respectively. Unknown EC digits are not
imputed. Cofactor-associated categories are recognized from reaction
participants and assigned confidence 0.4. They describe participant presence,
which is an imperfect proxy for a biochemical cofactor requirement.
Mechanism categories are coarse descriptors derived from atom-mapped bond
changes; mappings below confidence 0.5 are omitted. These descriptors do not
specify complete catalytic mechanisms. Cofactor and mechanism category
weights are $\xi_c=1$.

Reaction annotations are transferred to proteins only along training
associations. When several observations support the same membership, the
maximum confidence is retained. Native protein EC annotations can additionally
describe multiple functions of a training protein. Appendix~\ref{app:v4-annotations}
specifies the different annotation sources for the two benchmarks.

\paragraph{Relative within-category and background distances.}
Consider one endpoint type $X\in\{R,E\}$ and one family $f$. Let
$\hat{\mathbf{u}}_i=\vunit{\mathbf{u}_i^{\mathrm{tr}}}$ and
$\delta_{ij}=1-\hat{\mathbf{u}}_i^\top\hat{\mathbf{u}}_j$.
The index $i$ refers to a training endpoint of type $X$; these 512-dimensional
vectors differ from the complete retrieval vectors $\mathbf{z}_x$.
Let $\mathcal C_{X,f}$ contain the categories of family $f$ for endpoint type
$X$. Let $I_c$ index endpoints annotated with category $c$, $n_c=|I_c|$, and
$\mathcal A_f=\bigcup_{c\in\mathcal C_{X,f}}I_c$. The comparison set is
$B_c=\mathcal A_f\setminus I_c$. An endpoint without any annotation in $f$
is excluded from both membership and background for that family. For a
multiply annotated background endpoint, set
$\bar w_j^f=\max_{c'\in\mathcal C_{X,f}:j\in I_{c'}}w_{jc'}$.

For eligible categories, define the within-category and background confidence
masses
\begin{equation}
  T_c=\sum_{\substack{i,j\in I_c\\i\ne j}}w_{ic}w_{jc},\qquad
  U_c=\sum_{i\in I_c}\sum_{j\in B_c}w_{ic}\bar w_j^f,
\end{equation}
and confidence-weighted mean distances
\begin{align}
  d_c^{\mathrm{within}}
    &=\frac{1}{T_c}
       \sum_{\substack{i,j\in I_c\\i\ne j}}w_{ic}w_{jc}\delta_{ij},\\
  d_c^{\mathrm{background}}
    &=\frac{1}{U_c}\sum_{i\in I_c}\sum_{j\in B_c}
          w_{ic}\bar w_j^f\delta_{ij}.
  \label{eq:v4-bio-distances}
\end{align}
The per-category penalty is
\begin{equation}
  \ell_c=\bar w_c\,
       \vpos{m+d_c^{\mathrm{within}}-d_c^{\mathrm{background}}},
       \qquad m=0.1.
  \label{eq:v4-relative-bio}
\end{equation}
Here $\bar w_c=T_c/[n_c(n_c-1)]$ is the mean pairwise annotation confidence.
Its prefactor preserves the strength of annotation confidence. For example,
if every membership in a category has confidence 0.4, this prefactor is
$0.4^2=0.16$. Computing only confidence-normalized distances would remove
that common confidence reduction.

A category contributes only if $n_c\ge2$, $|B_c|\ge1$, $T_c>0$, and
$U_c>0$. Let $\mathcal C_{X,f}^{*}$ be the eligible categories. We use a
category-weighted average, followed by a fixed average over endpoint types:
\begin{equation}
  \mathcal L_{X,f}=
    \frac{\sum_{c\in\mathcal C_{X,f}^{*}}\xi_c\ell_c}
         {\sum_{c\in\mathcal C_{X,f}^{*}}\xi_c},\qquad
  \mathcal L_f=\tfrac12(\mathcal L_{R,f}+\mathcal L_{E,f}).
  \label{eq:v4-family-average}
\end{equation}
An empty eligible set contributes zero. Averaging by category, rather than
by the total number of pairs across categories, limits domination by very
large annotation groups. Membership overlap is allowed; the background for
one category can share other categories with its members.

\paragraph{Objective of the biological ablation.}
The ablation augments Eq.~\eqref{eq:v4-phase2-total} as
\begin{equation}
  \mathcal L_{\mathrm{Bio}}^{(2)}=\mathcal L^{(2)}
     +\frac13\sum_{f\in\mathcal F}\mathcal L_f.
  \label{eq:v4-phase2-bio-ablation}
\end{equation}
All three families have unit coefficients before the fixed average. This
recipe is recorded as \texttt{all\_1} in the implementation. In family-removal
ablation controls, the corresponding term is omitted; the divisor remains three.
Shuffled-label controls preserve annotation profiles, confidences, and
category frequencies while permuting their assignment to training endpoints.

The relative loss expresses a geometric preference for annotated peers. The
background does not assert that a protein is inactive for a reaction, and
different categories are not mutually exclusive biological functions.
Incomplete annotations and enzyme promiscuity can make these geometric
references imperfect. The relative hinge does not by itself guarantee
non-collapsed embeddings; association supervision remains necessary. In the
second training stage, biological gradients reach the residual heads but not
the frozen residue queries or feature extractors. Exact category-sum
identities avoid constructing a dense biological pairwise matrix
(Appendix~\ref{app:v4-efficient-bio}).

\subsection{Training-reaction dictionary representation}
\label{sec:v4-dictionary}
In addition to the neural embedding, both endpoints receive a vector
$\mathbf{a}\in\mathbb{R}^{N_R}$ whose coordinates index the target's training
reactions. The dictionary uses frozen raw features and observed training
associations. It is fitted once for each target split and is independent of
the current query or candidate collection.

\paragraph{Training-only centering.}
For a raw feature family with available training values $\mathbf{x}_i$, we compute
\begin{equation}
  \boldsymbol\mu=\frac{1}{N_{\mathrm{avail}}}\sum_{i\in\mathrm{available\ train}}\vunit{\mathbf{x}_i},
  \qquad \psi(\mathbf{x})=\vunit{\vunit{\mathbf{x}}-\boldsymbol\mu}.
  \label{eq:v4-center}
\end{equation}
The order is normalization, centering, then normalization. Enzyme features
are cached ProtT5 residue means. Reaction features concatenate four separately
centered and normalized blocks (ReactionT5v2, Uni-Mol2, ChIRo, chemistry),
followed by normalization of the concatenation. Unavailable blocks are zero.
For the dictionary's molecular blocks, the implemented export uses the
\emph{mean forward-reactant bank}: this is the complete participant set for
ReactZyme self-reactions and the reactant side for EnzymeMap. These fixed raw
features differ from the trainable side pooling in the neural reaction tower.

\paragraph{Enzyme-to-training-reaction transfer.}
Let $\psi_E(e)$ be the centered enzyme feature and let $\mathcal N_{32}(e)$
be its 32 nearest training proteins by cosine similarity
$c_{ej}=\psi_E(e)^\top\psi_E(e_j)$. With $\tau_E=0.03$, define
\begin{align}
  w_{ej}^{\mathrm{NN}}&=
    \exp\!\left(\frac{c_{ej}-\max_{t\in\mathcal N_{32}(e)}c_{et}}{\tau_E}\right),\\
  \tilde a_{e,i}&=
    \max_{\substack{j\in\mathcal N_{32}(e)\\Y^{\mathrm{tr}}_{ij}=1}}w_{ej}^{\mathrm{NN}},
  \qquad \mathbf{a}_e=\vunit{\tilde{\mathbf{a}}_e}.
  \label{eq:v4-enzyme-dictionary}
\end{align}
An empty maximum is zero. The maximum merges repeated support for a training
reaction without summing the frequency of similar supporting proteins.
There is no exclusion of an exact training protein when it is an eligible
neighbor. Such information is permitted by the benchmark's training graph.

\paragraph{Reaction-to-training-reaction response.}
Let $\psi_R(r)$ denote the normalized raw reaction feature and
$c_{ri}=\psi_R(r)^\top\psi_R(r_i)$. We use all training reaction anchors:
\begin{equation}
  \tilde a_{r,i}=\exp\!\left(
       \frac{c_{ri}-\max_t c_{rt}}{\tau_R}\right),\qquad
  \mathbf{a}_r=\vunit{\tilde{\mathbf{a}}_r},\qquad \tau_R=0.03.
  \label{eq:v4-reaction-dictionary}
\end{equation}
The implementation floors full-support exponential responses at the smallest
positive normal FP32 value before normalization. The rowwise maximum in
Eqs.~\eqref{eq:v4-enzyme-dictionary}--\eqref{eq:v4-reaction-dictionary}
stabilizes exponentiation and otherwise cancels under normalization.
Full reaction-anchor support avoids exact support mismatch caused by truncating
both endpoint vectors to separate nearest-anchor sets. It does not ensure
informative transfer for a reaction remote from the training chemistry.

\subsection{Inference adjustments and final bidirectional score}
\label{sec:v4-inference}
The reported recipe makes two explicit inference adjustments after both
training stages. First, it strengthens the entire fused protein branch:
\begin{equation}
  \tilde{\mathbf{b}}_e=\vunit{\mathbf{g}_e+\nu s \mathbf{f}_e},\qquad \nu=3,
  \qquad \tilde{\mathbf{b}}_r=\mathbf{b}_r.
  \label{eq:v4-inference-fusion}
\end{equation}
The scalar $\nu s$ is checked to remain below 0.5 for these checkpoints.
This multiplies the complete fused-view contribution, including its global,
SLEEC, and learned-slot information. It does not multiply only SLEEC weights.
Phase 2 was fitted to the native $\mathbf{b}_e$, before this adjustment.

Second, the existing residual heads are applied with inference multiplier
$c=0.5$:
\begin{equation}
  \mathbf{u}_e=\vunit{\tilde{\mathbf{b}}_e+c\kappa h_E(\tilde{\mathbf{b}}_e)},\qquad
  \mathbf{u}_r=\vunit{\mathbf{b}_r+c\kappa h_R(\mathbf{b}_r)},\qquad c\kappa=0.1.
  \label{eq:v4-inference-residual}
\end{equation}
The multiplier halves the residual coefficient relative to training. It is
not a hard bound on the update norm or on the resulting angular displacement.
The dictionary representation is unchanged by either adjustment.

With semantic weight $\alpha=0.4$, define independent endpoint embeddings
\begin{equation}
  \mathbf{z}_e=[\sqrt{1-\alpha}\,\mathbf{u}_e;\sqrt{\alpha}\,\mathbf{a}_e],\qquad
  \mathbf{z}_r=[\sqrt{1-\alpha}\,\mathbf{u}_r;\sqrt{\alpha}\,\mathbf{a}_r].
  \label{eq:v4-concatenation}
\end{equation}
Their dot product gives the final score
\begin{equation}
  \boxed{S(r,e)=\mathbf{z}_r^\top\mathbf{z}_e
       =0.6\,\mathbf{u}_r^\top \mathbf{u}_e+0.4\,\mathbf{a}_r^\top \mathbf{a}_e.}
  \label{eq:v4-final-score}
\end{equation}
Thus the neural block has 512 dimensions, whereas the complete representation
has $512+N_R$ dimensions. The same score is used in both retrieval directions.
Dense neural blocks and sparse enzyme dictionary blocks can be stored
separately without changing the mathematical score. The method requires the
frozen feature extractors, trained encoders, refinement heads, training
dictionary, and feature-normalization statistics at deployment.

\subsection{Optimization and target-specific fitting}
\label{sec:v4-optimization}
We fit a separate model for each ReactZyme split and for EnzymeMap, using the
corresponding target training associations. Trainable retrieval components
are freshly initialized for each target; the pretrained feature extractors
and SLEEC prior are reused and frozen. In particular, EnzymeMap does not
inherit a ReactZyme-trained F3 checkpoint.

Phase 1 uses seed 42, association batch size 512, AdamW learning rate
$10^{-4}$, weight decay $10^{-2}$, and constant learning rate. We use the
epoch-10 snapshot. The original ReactZyme trajectories continued to epoch 20;
their retained epoch-10 snapshots define this recipe. Features are cached,
training uses mixed BF16 precision, and contrastive calculations are kept in
FP32. Phase 2 uses seed 42, 100 full-graph AdamW updates, learning rate
$10^{-4}$, weight decay $10^{-3}$, and gradient-norm clipping at 1. The
declared step-100 state is used for the reported configuration.

The graph, dictionary centers, and dictionary associations use training
endpoints only. Added biological annotations are restricted to the separate
ablation. Matching target training associations does not imply identical
pretrained resources across methods. The present recipe was developed through exploratory experiments;
the fixed recipe should not be described as a configuration selected without
prior inspection of benchmark outcomes.
